% Week 3 section of the canonical synthesis exercise sheet.
% Included by Exercises.tex and by the lecture-specific wrapper.

\section{Week 3 --- Stochastic Integration: Construction and Applications}

\lecturesummary{Left-point gains and predictability; simple stochastic
integrals; the It\^o isometry and the \(L^2\) extension; the integrability
required for true martingales; conditional variance, stopping, and quadratic
variation; deterministic Gaussian integrals; deterministic integration by
parts; left/right/midpoint conventions; and simulation diagnostics. No
general It\^o formula is used.}

\subsection*{Part I --- Predictability and Simple Integrals}

\begin{exercise}[Level 1 --- Information available at the trading time]
Let \((W_t)_{t\geq0}\) be Brownian motion with its natural filtration. On the
grid \(0,1,2\), write \(\Delta W_i=W_{i+1}-W_i\).
\begin{exparts}
  \item For each rule below, determine whether the position used on
  \((i,i+1]\) is predictable:
  \begin{enumerate}[label=(\roman*),leftmargin=2.4em,itemsep=0.1em]
    \item \(\phi_0=1\), \(\phi_1=1+W_1^2\);
    \item \(\phi_0=\operatorname{sign}(W_1)\), \(\phi_1=0\);
    \item \(\phi_0=0\), \(\phi_1=\1_{\{W_1>0\}}\);
    \item \(\phi_i=\operatorname{sign}(\Delta W_i)\).
  \end{enumerate}
  Justify each answer with the relevant measurability statement.
  \item For every admissible rule, write the terminal gain explicitly as a
  weighted sum of Brownian increments.
  \item Suppose \(\phi_{t_i}\in\F_{t_i}\) and
  \(\E[\phi_{t_i}^2]<\infty\). Prove that
  \[
  \E[\phi_{t_i}\Delta W_i\mid\F_{t_i}]=0.
  \]
  \item Explain why the last identity prevents a predictable Brownian trading
  strategy from creating positive expected gain, but does not prevent it from
  changing risk.
\end{exparts}
\end{exercise}

\begin{exercise}[Level 1 --- Deterministic and random step strategies]
Consider the processes
\[
\phi_t=2\1_{(0,1]}(t)-\1_{(1,3]}(t)
\]
and
\[
\psi_t=\1_{\{W_1>0\}}\1_{(1,2]}(t).
\]
\begin{exparts}
  \item Verify that both processes are predictable.
  \item Write \(\int_0^3\phi_t\dd W_t\) and
  \(\int_0^2\psi_t\dd W_t\) without integral notation.
  \item Identify the law of the first integral.
  \item Compute the mean and second moment of the second integral.
  \item Show that the second integral has an atom at zero and identify its law
  as a mixture. Why is it not Gaussian?
  \item Explain why deterministic integrands behave differently from random
  predictable integrands with respect to Gaussianity.
\end{exparts}
\end{exercise}

\subsection*{Part II --- Isometry and the \(L^2\) Construction}

\begin{exercise}[Level 2 --- Prove the It\^o isometry]
Let
\[
\phi_t=\sum_{i=0}^{n-1}\xi_i\1_{(t_i,t_{i+1}]}(t),
\qquad
\xi_i\in\F_{t_i},
\qquad
\E[\xi_i^2]<\infty.
\]
Set \(I(\phi)=\sum_i\xi_i\Delta W_i\).
\begin{exparts}
  \item Prove that \(\E[I(\phi)]=0\).
  \item If \(i<j\), prove carefully that
  \[
  \E[\xi_i\Delta W_i\,\xi_j\Delta W_j]=0
  \]
  by conditioning on \(\F_{t_j}\).
  \item Prove the isometry
  \[
  \E[I(\phi)^2]=\E\int_0^T\phi_t^2\dd t.
  \]
  \item For another simple predictable process \(\psi\), derive
  \[
  \E[I(\phi)I(\psi)]=\E\int_0^T\phi_t\psi_t\dd t.
  \]
  \item If \(|\phi_t|\leq L\), prove the risk bound
  \[
  \operatorname{Var}(I(\phi))\leq L^2T.
  \]
\end{exparts}
\end{exercise}

\begin{exercise}[Level 2 --- Approximation and an exact \(L^2\) error]
Let \(f(t)=t\) on \([0,T]\), take the uniform grid \(t_i=iT/n\), and define
the left-step approximation
\[
f_n(t)=\sum_{i=0}^{n-1}t_i\1_{(t_i,t_{i+1}]}(t).
\]
\begin{exparts}
  \item Compute
  \[
  \int_0^T|f(t)-f_n(t)|^2\dd t
  \]
  exactly.
  \item Deduce from the isometry that
  \[
  \sum_{i=0}^{n-1}t_i\Delta W_i
  \longrightarrow \int_0^Tt\dd W_t
  \quad\text{in }L^2(\Omega).
  \]
  Give the exact root-mean-square error.
  \item Explain why the isometry makes the integral independent of the chosen
  sequence of simple predictable approximations.
  \item On \([0,1]\), compare
  \[
  \phi_t=1,
  \qquad
  \psi_t=\sqrt2\,\1_{(0,1/2]}(t).
  \]
  Show that the two terminal integrals have the same law but that their gain
  processes do not accumulate risk at the same times.
\end{exparts}
\end{exercise}

\subsection*{Part III --- Martingales, Integrability, and Stopping}

\begin{exercise}[Level 2 --- Audit the martingale assumptions]
Let \(\phi\) be predictable and suppose
\[
\E\int_0^T\phi_s^2\dd s<\infty.
\]
Define \(M_t=\int_0^t\phi_s\dd W_s\).
\begin{exparts}
  \item Use the isometry and Cauchy--Schwarz to prove that
  \[
  \E|M_t|
  \leq
  \left(\E\int_0^t\phi_s^2\dd s\right)^{1/2}<\infty.
  \]
  Explain why checking \(M_t\in L^1\) is logically necessary before using a
  conditional expectation in the definition of a martingale.
  \item Prove the martingale identity first for a simple predictable process,
  then explain how \(L^2\) approximation allows passage to a general \(\phi\).
  \item Show that \(L^2\) convergence implies the \(L^1\) convergence needed
  to pass conditional expectations to the limit.
  \item State clearly the difference between the assumptions
  \[
  \E\int_0^T\phi_s^2\dd s<\infty
  \qquad\text{and}\qquad
  \int_0^T\phi_s^2\dd s<\infty\quad\text{a.s.}
  \]
  What type of martingale is guaranteed by each condition?
\end{exparts}
\end{exercise}

\begin{exercise}[Level 3 --- Pathwise finite energy need not give a true martingale]
Let \(X\geq0\) be finite almost surely, independent of \(W\), and satisfy
\(\E[X]=\infty\). Enlarge the filtration so that \(X\in\F_0\), and set
\(\phi_t=X\).
\begin{exparts}
  \item Show that \(\phi\) is predictable and
  \[
  \int_0^T\phi_t^2\dd t=X^2T<\infty
  \quad\text{a.s.}
  \]
  \item Show formally, and then through localization, that
  \[
  M_t=\int_0^tX\dd W_s=XW_t.
  \]
  \item Prove that \(\E|M_t|=\infty\) for every \(t>0\). Deduce that \(M\)
  is not a true martingale.
  \item For \(n\geq1\), define
  \[
  \tau_n=
  \begin{cases}
  \infty,&X\leq n,\\
  0,&X>n.
  \end{cases}
  \]
  Verify that \(\tau_n\uparrow\infty\) almost surely and that
  \((M_{t\wedge\tau_n})_{t\geq0}\) is a square-integrable martingale.
  Conclude that \(M\) is a local martingale.
  \item Explain how the conclusions change if \(\E[X^2]<\infty\).
\end{exparts}
\end{exercise}

\begin{exercise}[Level 2 --- Conditional risk of a state-dependent strategy]
On \([0,2]\), define
\[
\phi_t=\1_{(0,1]}(t)
+\bigl(1+\1_{\{W_1>0\}}\bigr)\1_{(1,2]}(t),
\qquad
M_t=\int_0^t\phi_s\dd W_s.
\]
\begin{exparts}
  \item Verify predictability and expected finite energy.
  \item Compute \(M_1\) and write \(M_2-M_1\) explicitly.
  \item Compute \(\E[M_2\mid\F_1]\).
  \item Find \(\operatorname{Var}(M_2-M_1\mid\F_1)\) and the conditional law
  of \(M_2\) given \(\F_1\).
  \item Compute \([M]_2\) and interpret it as a state-dependent exposure
  clock.
\end{exparts}
\end{exercise}

\begin{exercise}[Level 2 --- Stop the exposure at a barrier]
Let \(\tau\) be a stopping time and define
\[
M_t=\int_0^t\1_{\{s\leq\tau\}}\dd W_s.
\]
\begin{exparts}
  \item Give the trading interpretation of the integrand and show that
  \(M_t=W_{t\wedge\tau}\).
  \item For fixed finite \(t\), verify
  \[
  \E\int_0^t\1_{\{s\leq\tau\}}^2\dd s
  =\E[t\wedge\tau]\leq t.
  \]
  Explain why this is the relevant integrability check.
  \item Deduce
  \[
  \E[W_{t\wedge\tau}]=0,
  \qquad
  \E[W_{t\wedge\tau}^2]=\E[t\wedge\tau].
  \]
  \item Apply the result to
  \(\tau_a=\inf\{u\geq0:W_u=a\}\), with \(a>0\), without using the
  density of \(\tau_a\).
  \item Explain why these conclusions do not require \(\tau\) itself to be
  bounded.
\end{exparts}
\end{exercise}

\subsection*{Part IV --- Deterministic Gaussian Applications}

\begin{exercise}[Level 2 --- Gaussian projection and independent risks]
Let
\[
I_T=\int_0^Tt\dd W_t.
\]
\begin{exparts}
  \item Identify the law of \(I_T\).
  \item Write \(W_T\) as a stochastic integral and compute
  \(\operatorname{Cov}(I_T,W_T)\).
  \item Find \(\E[I_T\mid W_T]\) and
  \(\operatorname{Var}(I_T\mid W_T)\).
  \item Find an affine function \(g(t)=t-c\) such that
  \[
  Y=\int_0^Tg(t)\dd W_t
  \]
  is independent of \(W_T\), and compute \(\operatorname{Var}(Y)\).
  \item Interpret the preceding calculations as orthogonal projection in
  \(L^2(0,T)\).
  \item For \(M_t=\int_0^ts\dd W_s\), show that
  \[
  (M_t)_{t\geq0}\overset{\mathrm{law}}=
  (B_{t^3/3})_{t\geq0}
  \]
  for a Brownian motion \(B\).
\end{exparts}
\end{exercise}

\begin{exercise}[Level 2 --- Deterministic integration by parts and Brownian area]
Let \(f:[0,T]\to\R\) be deterministic and absolutely continuous with
\(f'\in L^2(0,T)\).
\begin{exparts}
  \item Explain why \(f'\in L^1(0,T)\), why \(f\) has finite variation, and
  why \(f\in L^2(0,T)\).
  \item Starting from discrete summation by parts, prove
  \[
  f(T)W_T-f(0)W_0
  =\int_0^Tf(t)\dd W_t+\int_0^TW_tf'(t)\dd t.
  \]
  Make clear which limit is in \(L^2\) and which is pathwise.
  \item Explain why no general It\^o formula and no quadratic-variation
  correction are needed.
  \item With \(f(t)=t\), prove
  \[
  \int_0^TW_t\dd t
  =TW_T-\int_0^Tt\dd W_t.
  \]
  \item Identify the law of the Brownian area
  \(A_T=\int_0^TW_t\dd t\).
\end{exparts}
\end{exercise}

\subsection*{Part V --- Conventions and Simulation}

\begin{exercise}[Level 2 --- Left, right, and midpoint sums]
For a partition \(\pi\) of \([0,T]\), define
\[
L^\pi=\sum_iW_{t_i}\Delta W_i,
\quad
R^\pi=\sum_iW_{t_{i+1}}\Delta W_i,
\quad
S^\pi=\sum_i\frac{W_{t_i}+W_{t_{i+1}}}{2}\Delta W_i.
\]
\begin{exparts}
  \item Use
  \[
  W_{t_{i+1}}^2-W_{t_i}^2
  =2W_{t_i}\Delta W_i+(\Delta W_i)^2
  \]
  to identify the limit of \(L^\pi\).
  \item Deduce the limits of \(R^\pi\) and \(S^\pi\).
  \item Compute the mean of each limit. Which convention preserves the
  zero-mean martingale-gain interpretation?
  \item Let \(I_T=\int_0^TW_t\dd W_t\). Compute \(\E[I_T^2]\) from the
  explicit formula and again from the isometry.
  \item Explain exactly what is illegal about choosing
  \(\phi_{t_i}=\operatorname{sign}(\Delta W_i)\).
\end{exparts}
\end{exercise}

\begin{exercise}[Level 2 --- Simulation error and implementation checks]
On the uniform grid \(t_i=iT/n\), simulate
\[
\Delta W_i=\sqrt{\Delta t}\,Z_i,
\qquad
Z_i\overset{\mathrm{iid}}{\sim}\mathcal N(0,1),
\qquad
\Delta t=T/n.
\]
Define
\[
I^{(n)}=\sum_{i=0}^{n-1}W_{t_i}\Delta W_i,
\qquad
I=\int_0^TW_t\dd W_t.
\]
\begin{exparts}
  \item Prove the exact error identity
  \[
  I-I^{(n)}
  =\frac12\sum_{i=0}^{n-1}\bigl((\Delta W_i)^2-\Delta t\bigr).
  \]
  \item Derive
  \[
  \E[(I-I^{(n)})^2]=\frac{T^2}{2n}
  \qquad\text{and}\qquad
  \operatorname{RMSE}(I^{(n)})=\frac{T}{\sqrt{2n}}.
  \]
  \item For a state-dependent integrand \(\phi_t=h(W_t)\), write the correct
  order in which \(\phi_i\), \(I_{i+1}\), and \(W_{t_{i+1}}\) must be
  updated.
  \item Explain the effect of each implementation error:
  \begin{enumerate}[label=(\roman*),leftmargin=2.4em,itemsep=0.1em]
    \item using \(\Delta t\,Z_i\) instead of \(\sqrt{\Delta t}\,Z_i\);
    \item computing \(\phi_i\) from \(W_{t_{i+1}}\);
    \item using independent Brownian increments for the state and its
    stochastic integral;
    \item replacing \(\Delta W_i\) by \(\Delta t\).
  \end{enumerate}
  \item Propose at least four diagnostics involving means, second moments,
  covariance, grid refinement, or exact benchmarks. Explain why inspecting one
  plausible path is not sufficient.
\end{exparts}
\end{exercise}
